\section{Audit conclusions, integration changes, and further research}
\label{sec:audit-research}

\subsection{The precise change in mathematical status}

The inspection covered the source statements used by this investigation,
their nearby conjectures, and the six incoming reports at the pinned
revision. It was not a line-by-line certification of the entire collective
manuscript. No new counterexample or invalid proof in those inherited
statements was established. Several tempting shortcuts in the present
derivations were rejected or corrected before delivery. The appropriate
repository changes are the following additions and status refinements.

\begin{longtable}{@{}>{\raggedright\arraybackslash}p{.23\textwidth}>{\raggedright\arraybackslash}p{.34\textwidth}>{\raggedright\arraybackslash}p{.36\textwidth}@{}}
\toprule
Target in the inspected material & Proved advancement & Boundary retained\\
\midrule
\endhead
All reconstruction rows, Bernoulli report R2 &
Theorem~\ref{thm:all-rows} gives every entry of $M_b^{-1}$;
Theorems~\ref{thm:extraction-series} and \ref{thm:polylog-filter}
give its ordinary and polylogarithmic evaluations. &
The inverse is for the $(2,2b)$ matrix. Formal coordinate rank is
not arithmetic independence.\\[.5em]
Optimal-prefix conjecture beyond $p=6$ &
Theorem~\ref{rank8:theorem} proves the $p=8$, $q=2b$ family
for $b\ge5$, $b\not\equiv1\pmod {16}$. &
The missing residue class and general unequal orders $p\ge10$
are not settled by this paper.\\[.5em]
Centered cubic tails &
Theorem~\ref{cc:all} evaluates every regular odd centered cubic
trigamma moment; Theorems~\ref{cc:generator} and \ref{cc:polylog}
give its analytic and finite polylogarithmic generators. &
The lowest cubic evaluation has classical antecedents.
No general higher-power closure theorem is claimed.\\[.5em]
Next two-particle resonance &
Theorems~\ref{sr:thm:germ} and \ref{sr:thm:constant} establish
the grouped subtraction and exact cutoff constant. Corollary
\ref{sr:cor:alljets} gives every spectral row. &
The result concerns two particles and the specified next resonance.
Arbitrary particle number and arbitrary resonance remain open.\\[.5em]
Higher spatial spectral functionals &
Theorem~\ref{mellin:main} and Corollary~\ref{mellin:coefficients}
give ordinary logarithmic Mellin formulas at every coefficient order. &
The statement uses the inherited unshifted differential-polynomial
residue germ. General independent spatial shifts need a new analysis.\\[.5em]
Short Gaussian $S_6$ and revised $S_8$ &
No status change. &
The current frozen candidates remain conjectural. The older rejected
$S_8$ vector is a different statement.\\
\bottomrule
\end{longtable}

For integration, the full inverse belongs immediately after the rational
row in the Bernoulli moment-rank development. The rank-eight theorem is
an additional case of the existing optimal-prefix programme, and should
be attached with its congruence qualification. The cubic section belongs
after the canonical centered-square chapter. The resonance theorem
continues the determinantal report; the Mellin theorem continues the
differential-polygamma report. These are additive changes. The package's
integration notes identify the source files and supplied artifacts.

\subsection{Normalization issues that must survive editing}

\paragraph{The omitted binomial term.}
The lower-index-zero summand in the source tail matrix is not present.
Restoring it changes the matrix, the exceptional zeroth column, and the
reconstruction coefficients. Formula~\eqref{eq:M-system} and the direct
finite-cutoff verification in the package fix this convention explicitly.

\paragraph{The zeroth power in Faulhaber's formula.}
The uniformly correct expression for $\sum_{k=1}^N k^m$, $m\ge0$,
is \eqref{eq:faulhaber-tail}, with $B_{m+1}(1)$ at the lower endpoint.
Replacing it by $B_{m+1}(0)$ is valid for $m\ge1$ but fails at $m=0$.
This scope error was found in an early draft of the present article and
corrected. It is not reported as an error in an inherited theorem.

\paragraph{Exceptional cubic boundary terms.}
The $-\zeta(2)/2$ term at $j=2$ and the rational $1/4$ term at $j=3$
in \eqref{cc:Cj} are required. The latter is the surviving
$B_1^2/B_1^3$ boundary contribution in \eqref{cc:boundaryexception}.
Deleting low-order exceptions because a generic parity pattern looks
simpler would change the convergent identities.

\paragraph{Cutoff constants and spectral finite parts.}
At the second resonance, $C$ in \eqref{sr:eq:lambda-C} and $F$ in
\eqref{sr:eq:spectral-laurent} are different precisely normalized
quantities. Equation~\eqref{sr:eq:normalization} includes both derivatives
of the amplitude needed at a double pole. The example $a=c=1$ gives
the same cutoff constant $\gamma-\gamma^2$ but different spectral
finite parts under two normalizations. This is a proved comparison,
not an assertion that the inspected report made the incorrect identification.

\paragraph{A second-order endpoint constant after scalar cancellation.}
When $\phi_P=0$, the scalar-filter Stieltjes terms disappear, but
$\pi^2p_1(0)/6$ remains in \eqref{mellin:q2}. It comes from the sine
factor in the Mellin transform. Its absence at first order provides
no reason to discard it at second order.

\paragraph{An extrapolation that is actually false.}
The rank-eight determinant's valuation follows a simple law on the
proved residue classes, but the unrestricted extrapolation fails at
$b=65$ and $129$, as documented in the rank-eight section. These are
counterexamples to a guess considered during this investigation, not
to the source optimal-prefix conjecture. The determinants at those
points are nonzero. The missing class must be addressed by a different
argument or a finer congruence analysis.

\subsection{Exact certificates and independent diagnostics}

The accompanying scripts separate identities over $\Q$ from numerical
comparisons. The reconstruction verifier checks both products with the
original binomial matrix through $b=15$, the shifted and unshifted
orthogonality identities, the exceptional inverse row, and the printed
examples. A second implementation reconstructs 35 late-coordinate rows
directly from finite summation and the uncentered Hurwitz expansion;
it does not insert the proposed matrix formula. Exact coefficient tables
contain all 20 extraction rows for $2\le b\le6$.

The rank-eight computation is a finite symbolic certificate in an
integer polynomial ring, followed by a proved congruence for the Bernoulli
inputs. Checking all coefficients of that polynomial identity is what
makes it apply to an infinite family. A list of specialized nonzero
determinants would not suffice. The full polynomial and a generator that
reconstructs it from the four-by-four matrix are supplied.

The cubic verifier independently reconstructs 18 coordinates from finite
cutoffs, checks the generating-function coefficients, and differentiates
the finite polylogarithmic numerator twice with exact zero residual.
It also checks the rational and $\zeta(3)$ extraction identities.
The resonance verifier checks the finite Pearson relations and the
Laurent cancellations algebraically, and evaluates defining finite
weight sums with an independently generated asymptotic tail. The Mellin
verifier checks 56 exact principal-part and coefficient identities,
then compares spatial quadrature with spectral expressions.

Numerical diagnostics use 50, 80, or 110 decimal working digits, as
recorded by each script. The six reconstruction comparisons have largest
absolute residual below $3.4\times10^{-59}$; the cubic moment checks
through $r=8$ have largest higher-cutoff residual below
$4.8\times10^{-76}$. The two nonzero-parameter resonance comparisons,
including a complex parameter, agree within $1.4\times10^{-30}$.
The nine Mellin comparisons have largest relative residual below
$1.1\times10^{-49}$. These are independently organized diagnostics,
not rigorous interval enclosures. Cancellation can exhaust working
precision at large polynomial degree, so the reports retain the cutoff,
expansion order, and successive comparison rather than a single
unqualified count of ``correct digits.'' The analytic proofs establish
the infinite-family statements.

\subsection{Further research questions}

The following questions retain a concrete exact target. They prioritize
identities and reductions; estimates are tools for justifying the
manipulations, rather than separate objectives.

\begin{enumerate}[leftmargin=*,label=\textbf{Q\arabic*.}]

\item \textbf{Finish rank eight on $b\equiv1\pmod {16}$.}
Determine an arithmetic invariant of the four-by-four perturbation that
cannot vanish in this residue class. Higher modulus versions of the
Faulhaber congruence are available, but the counterexamples to the simple
valuation law show that its exponent cannot be predicted merely from
$v_2(b-1)$. A useful result would prove nonvanishing for all remaining
$b$, or identify a complete, exact valuation recursion with a termination
argument. Checking more values alone does not settle the question.

\item \textbf{Prove the optimal prefix for general unequal even orders.}
For $(p,q)=(2a,2b)$ with $a\ge5$, the inherited reduction leaves an
$a$-square perturbation determinant. Seek a structural factorization,
a congruence stable in $a$, or an algebraic interpretation of its
nonvanishing. The Bernoulli functional is indefinite, so a positive
moment-matrix argument cannot simply be imported. The result should
identify the exceptional arithmetic factors, not leave their nonvanishing
as an implicit assumption.

\item \textbf{Construct all higher-$p$ reconstruction rows.}
The kernel of Theorem~\ref{thm:kernel} and a finite-rank inverse formula
suggest a route for $(4,2b)$ and $(6,2b)$, whose rank is already known.
The target is an explicit coefficient formula comparable to
\eqref{eq:all-rows}, with no unknown minors. Translate those coefficients
into ordinary polynomial-subtracted polygamma sums with one prescribed
even-zeta coordinate. The independent finite-cutoff method in the package
can verify their normalization before an all-order proof is attempted.

\item \textbf{Extend the reconstruction to shifted and character tails.}
Replace $A_p(n)$ by a Hurwitz tail or a Dirichlet-character tail.
Determine the correct generalized Bernoulli kernel and state the exact
finite correction when the midpoint identity
$\zeta(-2j,1/2)=0$ is lost. One can first ask for a single rational
shift and a common character, then for distinct shifts. A claimed
ordinary evaluation must include the changed subtraction constants.

\item \textbf{Classify higher powers of the centered trigamma tail.}
For $d\ge4$, examine
$\sum_{n\ge0}\psi'(n+1)^d x_n^{-s}$ at its regular integer values.
Does the correct parity family lie in a fixed finite coordinate space
for each $d$? The cubic proof identifies a finite cancellation of linear
tails before taking constant terms. Find its higher-power analogue, or
prove a specific obstruction. An explicit quartic formula and a generator
would be a decisive next case; no quartic closure is assumed here.

\item \textbf{Determine the functional depth of the cubic generator.}
Theorem~\ref{cc:polylog} uses six specified depth-two polylogarithms.
Determine whether their particular combination admits a smaller function
basis under an explicitly stated class of argument transformations.
An equality of derivatives together with an anchored value is a suitable
certificate. A numerical relation at one $e^{-t}$ does not establish a
functional reduction, and formal symbol rank alone does not establish
arithmetic irreducibility at an endpoint.

\item \textbf{Integrate the new boundary operators at interior arguments.}
The operators in \eqref{eq:Phi-filter} prescribe exact cancellations at
$z=1$. Find anchored primitives that lower their differential order and
give compact formulas at interior algebraic points, including $1/2$ and
cyclotomic points accessible by continuation. Track all integration
constants from the power series at zero. The rational resonance generator
in \eqref{sr:eq:G-polylog} supplies a simple model of this approach.

\item \textbf{Develop a general simultaneous-escape subtraction.}
For three or more particles, classify the regions in which subsets of
summation indices tend to infinity together. Determine an exact
inclusion--exclusion subtraction and the associated exponents in the
Hurwitz terms. The doubled scale $1+2u$ in the present two-particle
case is the first nontrivial example. A successful theorem must prove
normal convergence of the full grouped expression before differentiating
at a resonance.

\item \textbf{Explain the saving of one shape-polygamma order.}
Equation~\eqref{sr:eq:lower-order-G} removes the apparent $R_{k+2}$
dependence from every $k$th spectral row. Is this a consequence of a
general determinant identity or of the two-particle Pearson reduction?
Find a corresponding regular expression at higher resonances or prove
where the saving fails. Such a result would reduce the number of
independent harmonic weights needed in exact derivative formulas.

\item \textbf{Extend the ordinary Mellin bridge to independent shifts.}
For $P$ in polygamma functions evaluated at several distinct shifts,
formulate the exact shift-difference principal parts and identify the
replacement for $D_j(u+j-1)$. A natural target is a normally convergent
Lerch or Hurwitz multiple-zeta family with all endpoint logarithms
explicit. The first problem is to specify the spatial subtraction scale
and its transformation law; the generalized Stieltjes corrections then
have an unambiguous meaning.

\item \textbf{Reduce the concrete nonlinear spectral derivatives.}
The cube identity \eqref{mellin:trigammacube} provides particular
combinations of $\partial_1^j Z(s,t)$ at convergent integer pairs.
Determine which combinations reduce to ordinary zeta derivatives,
height-one coefficients, or the existing Tornheim coordinates. Functional
relations must hold in a neighborhood of the differentiated point;
an identity true only at integer orders cannot be differentiated as if
it were such a relation. The logarithmic integral
\eqref{mellin:cube-log} is a concrete initial target.

\item \textbf{Compute the spectral image of the differential quotient.}
The incoming report classifies integration by parts in a free
differential-polynomial algebra. The Mellin bridge supplies a map from
that quotient, with its endpoint functional, to harmonic residue series.
Determine the kernel of this map at fixed weight, first exactly at low
weight and then by structural generators. A formal quotient dimension
and a dimension after substituting the digamma function need not agree.

\item \textbf{Build anchored multivariable primitive systems.}
Combine the spectral derivatives in \eqref{sr:eq:alljets} with the
shift primitive \eqref{sr:eq:primitive}. Determine finite recurrences for
iterated primitives in the shift and shape parameters, with base values
fixed at specified anchors. The goal is an identity in known Gamma,
Hurwitz-zeta, Stieltjes, and multiple-zeta families, rather than naming
an undetermined constant after each integration.

\item \textbf{Return to the frozen Gaussian $S_6$ and revised $S_8$ targets.}
Preserve their exact current vectors and the repository's established
shuffle equivalences. Seek a functional or iterated-integral relation
whose rigorously evaluated boundary is precisely one of those vectors.
The present identities expand the available analytic tools but establish
no direct implication for those two conjectures. Their status should
change only after such a relation and its endpoint evaluation are proved.

\end{enumerate}

The concrete deliverable of this investigation is therefore a set of
proved reconstruction, moment, resonance, and Mellin theorems with
explicit coefficients and normalization constants. The unresolved
questions are kept attached to the exact missing mathematical step.
